% Shared section of the physics documents of the six fixed-B modes: derivatives of the result (the adjoint).
% \input by each mode document after \adjointcolumntrue (column modes) or \adjointpoissontrue (Poisson modes).
\section{Derivatives of the result: the adjoint}\label{sec:adjoint}
An inversion that fits observed spectra changes $T$ (and $\B$) and asks how a misfit $\ell$ that depends on the atmosphere,
through $P$, $\rho$, $n_e$, $\chi_{500}$, $\tau_{500}$ and the heights $z$, changes with them. It needs the gradient
$\partial\ell/\partial\ln T$ at every grid point, often $10^6$ numbers. Perturbing each input in turn and solving again would
cost one or two solves per number. The \textbf{adjoint} gives all of them for the cost of about one solve.

\subsection{The idea}
Write the solve as a map from the inputs $x=(\ln T,\B)$ to the outputs $y$, and the misfit as $\ell(y)$. By the chain rule
\begin{equation}
\frac{\partial\ell}{\partial x}=\Big(\frac{\partial y}{\partial x}\Big)^{\!\top}\bar y,\qquad \bar y\equiv\frac{\partial\ell}{\partial y}.
\label{eq:adj}
\end{equation}
$\bar y$ (the \emph{cotangent} of the outputs) is given by the caller. The Jacobian $\partial y/\partial x$ is never formed:
the transposed product is evaluated by running the computation backwards, each step passing the cotangent of its outputs
to its inputs. This is reverse-mode differentiation, and its cost does not depend on how many inputs there are.

\subsection{What is differentiated}
\ifadjointcolumn
The column integration is a fixed sequence of arithmetic steps (substeps, table look-ups, the top closure, the alignment and
the gauge), so its adjoint is the derivative of exactly that computation: the \textbf{discrete adjoint}. Its discrete choices
are recorded during the solve and held fixed: the number of substeps of each interval, the table cell of each look-up, the
branch of the $\tau$ top closure, the alignment level and its median column, and the interpolation brackets. Away from the
points where one of these choices switches, the map is smooth and this is its exact derivative; at a switch no derivative
exists. Where a step length depends on the state (the Heun substeps of the $\tau$ and mass axes), the length is
differentiated too. An implicitly defined quantity (the hydrostatic top closure) is differentiated by the implicit function
theorem. The equation-of-state derivatives are those of the table interpolant (its bicubic Hermite polynomials), not of the model behind the table.
\fi
\ifadjointpoisson
The Poisson solve is iterative: its result is the state $v^*$ at which the gradient of $\Phi$ vanishes,
$G(v^*,x)\equiv\partial\Phi/\partial v=0$. Differentiating this condition (the \textbf{implicit function theorem}) gives
\begin{equation}
\frac{\partial v^*}{\partial x}=-\mathsf H^{-1}\,\frac{\partial G}{\partial x},\qquad
\mathsf H=\frac{\partial^2\Phi}{\partial v^2},\label{eq:ift}
\end{equation}
so the adjoint solves one linear system with the Hessian $\mathsf H$, $\mathsf H\lambda=\bar v$, and returns
$-(\partial G/\partial x)^{\top}\lambda$, plus the contributions of everything else that depends on $x$: the starting
column-hydrostatic state (which sets the weights of $\Phi$ through their layer medians), the part above the cut and the
gauge. The iterations themselves (damping, step scaling, preconditioners, warm starts) only steer the solver toward $v^*$
and are not differentiated. Two points matter:
\begin{itemize}
\item \textbf{The full Hessian.} With $\B$ fixed the residual does not vanish at the solution (the Lorentz force is in general
not a gradient), so the Gauss--Newton matrix $\mathsf J^{\top}\mathsf W\mathsf J$ of the solver is not the Hessian: the second
derivatives of the residuals contribute too. The adjoint uses the full $\mathsf H$, applied to vectors by differences of the
exact gradient, and solves for $\lambda$ by preconditioned conjugate gradients.
\item \textbf{Directions $\Phi$ does not see.} Some changes of the unknowns leave $\Phi$ unchanged and $\mathsf H$ singular along
them. They keep their starting values, so the derivative passes through to the start along them: on the height grid the mean
$\ln P$ of the anchor plane; on the $\tau$ and mass axes a uniform shift of all heights (fixed by the gauge) and the shortest
horizontal waves of the centred differences, which the differences cannot see.
\end{itemize}
Equation~\eqref{eq:ift} holds only at a stationary point. The Gauss--Newton iteration damps every step and measures its own
progress with that damping, so a slowly converging direction can stay far from stationary while the iteration reports
little left to gain. The option \code{polish} therefore adds up to $n$ \textbf{Newton steps with the exact Hessian} after the
solve, which bring the returned state to $v^*$ (on a small test cube four steps lower the undamped Newton decrement to
$10^{-25}$). The adjoint reports how stationary the returned state was: \code{:stationary} when one undamped Newton step
would lower $\Phi$ by at most a relative $10^{-6}$, \code{:not\_stationary} otherwise, in which case the result is the
derivative of the point the solver was approaching.
\fi

\subsection{How exact it is}
The adjoint is checked against finite differences of the solve itself: $\langle\bar y,\,[y(x+h\,\delta x)-y(x-h\,\delta x)]/2h\rangle$
against $\langle\partial\ell/\partial x,\,\delta x\rangle$ for random $\delta x$ and $\bar y$, on 1, 2 and 3 MPI ranks and every
viewing geometry.
\ifadjointcolumn
The two agree to $3\times10^{-10}$--$3\times10^{-8}$ (relative) on all three axes.
\fi
\ifadjointpoisson
With \code{polish} they agree to $10^{-9}$--$10^{-7}$ (relative).
\fi
A solve with a given electron density (no hydrostatic integration) is forward only and has no adjoint.
